\phantomsection
\chapter*{Résumé et indexation en français}\label{chap:Résumé}
\addcontentsline{toc}{chapter}{Résumé et indexation en français}

Accord hybride des genres en polonais

Résumé : 

L'objectif de cette thèse est d'étudier l'accord hybride, tel que décrit par Greville Corbett dans ses travaux, en mettant l'accent sur le genre en polonais \citep{corbett_agreement_2023, corbett_agreement_1979}. Nous avons mené une étude de corpus en utilisant trois types de noms en polonais pouvant présenter un accord de genre hybride : les noms à genre social, les noms de profession et les formes péjoratives. Sur la base de 1 722 occurrences d'accord avec ces noms trouvées dans le Corpus national de polonais, nous avons constaté que la majorité des occurrences présentaient un accord syntaxique (74 \%). Nous avons également constaté des effets de l'ordre cible-contrôleur sur le choix entre l'accord sémantique et l'accord syntaxique dans l'ensemble des données. Nous avons également constaté un effet de la distance dans les occurrences d'accord avec les noms à genre social et un effet du type de cible dans les occurrences avec les noms de profession. Nous avons ensuite mené deux expériences de jugement d'acceptabilité, l'une avec des noms de profession et l'autre avec des formes péjoratives. Contrairement à l'étude de corpus, nous avons constaté une préférence générale pour l'accord sémantique. Ces études montrent que l'accord sémantique est acceptable en polonais et sensible à plusieurs facteurs de la hiérarchie d'accord de Corbett. Les résultats fournissent également des pistes pour des études futures sur le thème de l'accord de genre hybride en polonais.
\\

Mots clés français : accord ; genre ; polonais ; corpus ; expériences
\\

Forme ou Genre : thèse
\\

fMeSH\footnote{Répertoire structuré de mots-clés pour l'analyse de contenu et le classement de documents dans le domaine biomédical (français/anglais)} : Dissertation universitaire
\\

Rameau\footnote{Répertoire d'autorité-matière encyclopédique et alphabétique unifié (français)} :  Thèses et écrits académiques 	










